\PassOptionsToPackage{hyphens}{url}
\documentclass[aps,prd,twocolumn,superscriptaddress,nofootinbib]{revtex4-2}

% MiKTeX REVTeX 4.2 needs explicit package loading
\usepackage{amsmath}
\usepackage{amssymb}
\usepackage{amsthm}
\usepackage{booktabs}
\usepackage{tabularx}
\usepackage{xcolor}
\usepackage{graphicx}
\graphicspath{{../figures/}}

\newtheorem{definition}{Definition}
\newtheorem{proposition}{Proposition}
\newtheorem{theorem}{Theorem}
\newtheorem{corollary}{Corollary}
\newtheorem{lemma}{Lemma}
\newtheorem{remark}{Remark}
\newtheorem{axiom}{Axiom}

\usepackage{hyperref}
\hypersetup{
    pdftitle={The Saturation Theorem: Axiomatic Constraints on Quantum Gravity from Cosmological Relaxation},
    pdfauthor={Lukas Geiger},
    colorlinks=true,
    linkcolor=black,
    urlcolor=blue,
    citecolor=black
}

\begin{document}

% ===================================================================
% TITLE PAGE
% ===================================================================

\title{The Saturation Theorem: Axiomatic Constraints on Quantum Gravity\\from Cosmological Relaxation}

\author{Lukas Geiger}
\thanks{ORCID: \url{https://orcid.org/0009-0005-7296-1534}}
\affiliation{Independent Researcher, Bernau, Germany}

\date{\today}

\begin{abstract}
We formulate four minimal axioms -- finite geometric capacity, cooperative reinforcement, monotone cosmic control, and renormalization-group composition -- that any quantum gravity theory must satisfy if its macroscopic limit reproduces general relativity with a bounded relaxation mechanism. From these axioms alone, we derive the \textit{Saturation Theorem}: the unique smooth, scale-coherent effective response function whose composition extends regularly to the saturation boundary is the hyperbolic tangent $S(g) = S_{\max}\tanh(\alpha\,u(g))$, up to reparametrization. Alternative profiles -- logistic, rational, arctangent, or hard-saturation -- are excluded unless they reduce to $\tanh$ via admissible reparametrization. The key mathematical lever is a composition law (Axiom~D) restricting the functional form to solutions of Abel's equation, whose unique bounded monotone solution is $\tanh$. Major quantum gravity programs -- Loop Quantum Gravity, asymptotic safety, string/holographic approaches, causal set theory, and noncommutative geometry -- naturally satisfy the axioms when restricted to cosmological observables, explaining the universality of saturation dynamics across disparate frameworks. The result elevates the $\tanh$ saturation ODE $dX/da = k(1-X^2)$ of the Curvature Relaxation Model from a phenomenological fit to a \textit{universal normal form} of capacity-limited cooperative relaxation, the open question being which UV completion selects $k$ and $\Phi_0$.
\end{abstract}

\keywords{quantum gravity, saturation dynamics, axiomatic cosmology, No-Go theorem, Curvature Relaxation Model, tanh universality, renormalization group}

\thanks{AI tools were used for mathematical formalization and text generation; see Acknowledgments for details. All scientific content is the sole responsibility of the author.}

\thanks{Paper~V in the CRM series. Companion papers: \cite{Geiger2026,Geiger2026b,Geiger2026c,Geiger2026d}.}

\maketitle


% ===================================================================
% 1. INTRODUCTION
% ===================================================================
\section{Introduction}
\label{sec:intro}

\subsection{The problem}

The unification of quantum mechanics and general relativity remains the central open problem of theoretical physics. After decades of effort, several major programs have produced deep structural insights without converging on a single theory:

\begin{itemize}
\item \textbf{Effective field theory (EFT) of gravity} computes reliable quantum corrections at low energies but provides no UV completion \cite{Donoghue1994}.
\item \textbf{Asymptotic safety} seeks a UV fixed point of the gravitational renormalization group (RG) flow, constraining the theory without new degrees of freedom \cite{Weinberg1979,Reuter1998,Percacci2017}.
\item \textbf{String theory} provides a UV-finite framework containing gravity automatically, but faces the landscape/selection problem \cite{Polchinski1998}.
\item \textbf{Loop Quantum Gravity (LQG)} achieves background independence with discrete geometry operators, but struggles with dynamics and the classical limit \cite{Rovelli2004,Thiemann2007}.
\item \textbf{Holography/AdS-CFT} offers a nonperturbative definition of quantum gravity in anti-de~Sitter space, with unclear extension to realistic cosmology \cite{Maldacena1999}.
\item \textbf{Causal set theory} posits discrete causal structure as fundamental, with the continuum as emergent \cite{Sorkin2003,Surya2019}.
\item \textbf{Noncommutative geometry} reformulates geometry as spectral data, naturally bridging operators and spacetime \cite{Connes1994}.
\end{itemize}

Orthogonal to these UV-completion programs, the \textit{Swampland conjectures} \cite{Obied2018} constrain which low-energy theories can be consistently embedded in quantum gravity, while the \textit{EFT of Dark Energy} \cite{Gubitosi2013} parametrizes deviations from $\Lambda$CDM at the perturbative level. Neither framework, however, determines the \textit{functional form} of the cosmological saturation profile.

A striking observation is that when any of the UV-completion programs listed above is applied to cosmology, the resulting effective dynamics generically exhibits \textit{saturation behavior}: a bounded, cooperative approach to a maximal geometric configuration. This paper asks whether this universality is accidental or inevitable.

\subsection{The CRM as empirical anchor}

The Curvature Relaxation Model (CRM) \cite{Geiger2026,Geiger2026b,Geiger2026c,Geiger2026d} provides a concrete, empirically tested cosmological framework in which saturation dynamics plays the central role. The model replaces the cosmological constant $\Lambda$ and particulate dark matter with geometric curvature return, achieving $\Delta\chi^2 = -5.5$ vs.\ $\Lambda$CDM on joint SN+CMB+BAO data with 6 parameters and zero Early Dark Energy \cite{Geiger2026b}. The entire framework derives from a single dynamical equation:
\begin{equation}
\frac{dX}{da} = k(1 - X^2)\,,
\label{eq:saturation_ode}
\end{equation}
where $X = \Omega_\Phi/\Phi_0$ is the normalized curvature return potential and $a$ is the scale factor. The solution $X(a) = \tanh(k(a - a_{\mathrm{trans}}))$ provides late-time acceleration (``dark energy'') in its saturated phase and gravitational scaffolding (``dark matter'') in its unsaturated phase \cite{Geiger2026b}.

Paper~III \cite{Geiger2026c} derives this ODE from the effective Lagrangian $\mathcal{L}_{\mathrm{CRM}} = R/(16\pi G) + \gamma \mathcal{F}(T/\rho) R^2 + \frac{1}{2}(\partial\phi)^2 - V_{\mathrm{PT}}(\phi)$, where $V_{\mathrm{PT}}(\phi) = V_0/\cosh^2(\phi/\phi_0)$ is a P\"oschl-Teller potential. Paper~III also surveys five UV-completion candidates (LQG, Finsler geometry, information-theoretic spacetime, causal sets, quantum error correction), observing that all produce ODE structures of the form $dX/dt \propto (1 - X^2)$.

The present paper elevates this observation from \textit{structural analogy} to \textit{mathematical theorem}: we prove that any theory satisfying four physically motivated axioms must produce $\tanh$-type saturation, regardless of its microscopic mechanism.

\subsection{Scope and structure}

The central result is the \textit{Saturation Theorem} (Theorem~\ref{thm:saturation}): Under Axioms~A--D (including boundary regularity of the composition law), the effective cosmological response function is uniquely determined to be $\tanh$, up to reparametrization. The paper is organized as follows:

\begin{itemize}
\item Section~\ref{sec:axioms}: Formulation of the four axioms.
\item Section~\ref{sec:theorem}: Statement and proof of the Saturation Theorem.
\item Section~\ref{sec:exclusion}: Exclusion of alternative saturation profiles.
\item Section~\ref{sec:qg}: Verification of the axioms in major QG programs.
\item Section~\ref{sec:ontology}: Ontological implications: spacetime as an ordered phase.
\item Section~\ref{sec:observational}: Observational consequences and open questions.
\item Section~\ref{sec:discussion}: Discussion and outlook.
\end{itemize}


% ===================================================================
% 2. THE FOUR AXIOMS
% ===================================================================
\section{The Four Axioms}
\label{sec:axioms}

We formulate the axioms in terms of an effective response function $S_T(g)$, where $T$ denotes a quantum gravity theory and $g$ is an effective coupling/curvature/backreaction parameter measuring the ``distance'' from the UV to the IR. The function $S_T(g)$ encodes how the theory's geometric output (curvature, expansion rate, geometric order) responds to $g$.

\subsection{Axiom A: Finite Geometric Capacity}

\begin{axiom}[Finite Geometric Capacity]
\label{ax:capacity}
Per Hubble patch, there exists a finite maximum geometric order $C_{\max} < \infty$. Equivalently, the effective response function $S_T(g)$ is globally bounded:
\begin{equation}
0 \leq S_T(g) < S_{\max} < \infty \quad \forall\, g \geq 0\,.
\end{equation}
\end{axiom}

\textit{Physical content.} This axiom excludes infinite curvature, infinite information density, and UV divergences. It is extremely weak: it holds in LQG (finite area quanta \cite{Rovelli2004}), in holography (Bekenstein bound \cite{Bekenstein1981}), in quantum error-correcting spacetime (finite code capacity \cite{Almheiri2015}), in causal set theory (finite density \cite{Sorkin2003}), and in asymptotic safety (finite effective action at the UV fixed point \cite{Reuter1998}).

Without finite capacity, quantum gravity is by definition unstable: unbounded geometric configurations correspond to singularities, which every QG program aims to resolve.

\textit{Connection to CRM.} The saturation amplitude $\Phi_0$ in the CRM \cite{Geiger2026} is the physical realization of $S_{\max}$. The P\"oschl-Teller potential $V(\phi) = V_0/\cosh^2(\phi/\phi_0)$ enforces $\Omega_\Phi \leq \Phi_0$ dynamically \cite{Geiger2026c}.

\subsection{Axiom B: Cooperative Reinforcement}

\begin{axiom}[Cooperative Reinforcement]
\label{ax:cooperative}
Existing geometric order facilitates further ordering. The effective response function $S_T(g)$ is smooth ($C^\infty$) for $g \geq 0$, with a local expansion at $g = 0$:
\begin{equation}
S_T(g) = c_1 g + c_3 g^3 + \mathcal{O}(g^5)\,, \quad c_1 > 0\,.
\label{eq:ir_expansion}
\end{equation}
The odd-power structure reflects the catalytic nature of the response: the rate of geometric ordering is proportional to both the existing ordered fraction and the remaining capacity.
\end{axiom}

\textit{Physical content.} This axiom encodes the cooperative, self-reinforcing character of geometric ordering. It is the gravitational analog of ferromagnetic ordering, where aligned spins facilitate further alignment. Formally, the condition $c_1 > 0$ ensures that the response is positive (ordering, not disordering) for small perturbations, while the odd-power structure is consistent with the antisymmetric extension (Axiom~B').

Without cooperative reinforcement, there would be no classical spacetime: a linear or anti-cooperative response cannot produce the long-lived, large-scale geometric structures observed in the universe. The cooperative character is analogous to spontaneous symmetry breaking in condensed matter: local ordering seeds further ordering, driving the system toward a globally ordered state (saturation).

\textit{Connection to CRM.} The saturation coupling $k$ in the CRM ODE~\eqref{eq:saturation_ode} is $c_1$. The game-theoretic equilibrium of Paper~I \cite{Geiger2026} provides the thermodynamic motivation: the null space's concentration gradient $G$ drives cooperative ordering until the spacetime bubble saturates.

A companion condition extends the domain of $S_T$ to negative arguments, enabling the group-theoretic classification of Step~2 in the proof:

\begin{axiom}[Antisymmetric Extension (B')]
\label{ax:antisymmetry}
The response function $S_T$ admits a smooth antisymmetric extension to all $g \in \mathbb{R}$:
\begin{equation}
S_T(-g) = -S_T(g) \quad \forall\, g\,.
\label{eq:antisymmetry}
\end{equation}
\end{axiom}

\textit{Physical content.} The antisymmetry has three complementary motivations:

\textit{(i)~Sign covariance of the response.} The parameter $g$ is not a ``distance'' (which would be non-negative by definition) but a \textit{signed deviation} from a symmetric reference state. In the CRM, $g$ parametrizes the curvature perturbation relative to the de~Sitter background; positive $g$ corresponds to over-dense perturbations that enhance geometric ordering, and negative $g$ to under-dense regions that reduce it. Sign covariance ($S_T(-g) = -S_T(g)$) then states that the response function treats over- and under-dense regions symmetrically in magnitude -- a natural physical requirement absent any symmetry-breaking mechanism in the response itself. We emphasize that Axiom~B' is \textit{not} a claim about cosmological time-reversal symmetry (which is broken by the expansion); it is a constraint on the functional form of the response under $g \to -g$.

\textit{(ii)~Precedents in physics.} Sign covariance of response functions is ubiquitous: in electrodynamics, the polarization $P(-E) = -P(E)$ for centrosymmetric media; in magnetism, $M(-H) = -M(H)$ for the paramagnetic/ferromagnetic response; in the renormalization group, beta functions of dimensionless couplings are odd functions of the coupling deviation from a fixed point. The antisymmetry of Axiom~B' is the gravitational analog of these well-established symmetries.

\textit{(iii)~Mathematical role.} Together with the odd-power structure of Axiom~B, B' extends the composition law from the half-line $[0, S_{\max})$ to a full group operation on $(-S_{\max}, S_{\max})$, which is essential for the Abel-equation classification. Without B', only a semigroup structure is available, and the uniqueness argument of Step~2 does not apply in the same form. We note that this is an explicit assumption, not a derived consequence; its relaxation leads to alternative composition structures (e.g., $C(x,y) = x + y - xy$ on $[0,1)$), which is why it appears as a separate axiom rather than being subsumed into Axiom~B.

\subsection{Axiom C: Monotone Cosmic Control}

\begin{axiom}[Monotone Cosmic Control]
\label{ax:monotone}
The effective response function $S_T(g)$ is monotonically increasing, with saturation at infinity:
\begin{equation}
S_T'(g) \geq 0\,, \quad \lim_{g \to \infty} S_T(g) = S_{\max}\,.
\label{eq:monotone}
\end{equation}
There exists exactly one relevant direction (control parameter) governing the transition from low to high geometric order.
\end{axiom}

\textit{Physical content.} This excludes oscillating, multi-fixed-point, or chaotic geometric dynamics. It guarantees:
\begin{itemize}
\item A unique time direction (thermodynamic arrow).
\item No ``bouncing'' between geometric phases.
\item Stable expansion toward a de~Sitter attractor.
\end{itemize}

In the CRM, the control parameter is the scale factor $a$; in RG language, it corresponds to the single relevant direction flowing from the UV fixed point to the IR \cite{Percacci2017}. Axiom~C is violated by cyclic cosmologies and by models with multiple stable vacua -- both of which lack a unique thermodynamic arrow. We note that Axiom~C is a restriction to the \textit{cosmological sector} with a single dominant transition: it does not address the string landscape (where tunneling between metastable vacua produces non-monotone trajectories) or stochastic inflation (where local fluctuations break monotonicity). These scenarios involve multiple effective control parameters or stochastic dynamics, placing them outside the scope of the single-field, deterministic axiom system. The theorem should be read as: ``within the class of theories with a single deterministic saturation channel, the functional form is uniquely $\tanh$.''

\textit{Connection to CRM.} The monotonicity of $\Omega_\Phi(a)$ is a direct consequence of the ODE~\eqref{eq:saturation_ode}: since $1 - X^2 > 0$ for $|X| < 1$, the solution is strictly increasing.

\begin{remark}[Axiom interdependence]
As shown in Corollary~\ref{cor:monotonicity}, monotonicity (Axiom~C) follows from Axioms~A, B, B', and~D. We nevertheless retain Axiom~C as an explicit statement for physical transparency: it connects directly to the thermodynamic arrow of time and makes the exclusion of oscillating/chaotic cosmologies visible at the axiom level. The logical content of the axiom system is thus: three independent axioms (A, B/B', D) plus one derived consequence (C).
\end{remark}

\subsection{Axiom D: Renormalization-Group Composition}

\begin{axiom}[Renormalization-Group Composition]
\label{ax:composition}
There exists a scale composition $g \mapsto g_1 \oplus g_2$ (``two RG steps in succession'') such that the response function is compositionally coherent:
\begin{equation}
S_T(g_1 \oplus g_2) = \mathcal{C}\!\left(S_T(g_1),\, S_T(g_2)\right)\,,
\label{eq:composition}
\end{equation}
where $\mathcal{C}$ is a smooth, commutative, associative composition with neutral element $0$ and absorbing boundary $S_{\max}$:
\begin{equation}
\mathcal{C}(S, 0) = S\,, \quad \mathcal{C}(S, S_{\max}) = S_{\max}\,.
\label{eq:composition_boundary}
\end{equation}
Furthermore, $\mathcal{C}$ extends smoothly (i.e., $C^\infty$) to the closed square $[0, S_{\max}]^2$.
\end{axiom}

\textit{Physical content.} This is the deepest axiom. It states that the geometric response must be \textit{self-consistent under coarse-graining}: observing the universe at two successive scales and combining the results must give the same answer as observing at the combined scale directly. This is the defining property of a renormalization group.

The composition law~\eqref{eq:composition} encodes the requirement that the theory be \textit{scale-coherent} in the following precise sense: the response function $S_T$ is a homomorphism from the additive structure of the control parameter to the composition algebra of responses; i.e., no preferred scale exists at which this homomorphism property breaks down or changes character. The absorbing boundary condition $\mathcal{C}(S, S_{\max}) = S_{\max}$ reflects the fact that fully saturated geometry cannot be further ordered -- it is a fixed point.

The $C^\infty$-extension requirement has a direct physical interpretation: it demands that the composition law remains well-behaved \textit{even near saturation}. In physical terms, this means that the approach to the geometric capacity limit is smooth and does not introduce new divergences or critical exponents. Three physical arguments support this requirement:

\textit{(a)~No new phase transition at saturation.} A composition that develops singular derivatives at the boundary would signal a phase transition \textit{within} the saturation process itself -- i.e., the saturation mechanism would break down precisely where it is needed most. The smoothness condition excludes such pathological behavior. This is distinct from phase transitions at the \textit{onset} of saturation (which are permitted and correspond to the CRM's transition epoch $a_{\mathrm{trans}}$); the boundary regularity concerns the \textit{fully saturated} regime.

\textit{(b)~Distinction from critical phenomena.} While critical exponents and divergent susceptibilities characterize phase transitions at finite temperature (the Curie point in ferromagnets, $T_c$ in superfluids), the saturation boundary $S \to S_{\max}$ is an \textit{asymptotic} regime ($g \to \infty$), not a critical point. In the ferromagnetic analogy, it corresponds to $T \to 0$ (complete alignment), not to $T \to T_c$ (critical fluctuations). The zero-temperature limit of the magnetization \textit{is} smooth ($M \to M_s$ exponentially in mean-field theory and as a power law in low $T$, but without divergent derivatives of the \textit{composition} of magnetizations). The $C^\infty$ requirement is therefore physically appropriate for the saturation boundary.

\textit{(c)~Analyticity of the vacuum state.} In quantum field theory, the vacuum is an analytic state: all correlation functions at the vacuum are smooth. Since full saturation ($S = S_{\max}$) represents the geometric analog of the vacuum (the maximally ordered state), the composition of responses near this state should inherit the smoothness of the vacuum.

\textit{Connection to CRM.} In the CRM, the composition law emerges from the thermodynamic structure: the curvature return potential $\Omega_\Phi$ at scale $a_1 \cdot a_2$ is determined by the potentials at $a_1$ and $a_2$ through the saturation ODE's group property. The $\tanh$ function satisfies $\tanh(u_1 + u_2) = (\tanh u_1 + \tanh u_2)/(1 + \tanh u_1 \tanh u_2)$, which is precisely a composition law of the required form with $\mathcal{C}(x,y) = (x + y)/(1 + xy/S_{\max}^2)$.


% ===================================================================
% 3. THE SATURATION THEOREM
% ===================================================================
\section{The Saturation Theorem}
\label{sec:theorem}

\subsection{Statement}

\begin{theorem}[Saturation Theorem]
\label{thm:saturation}
Let $T$ be a theory whose effective response function $S_T(g)$ satisfies Axioms~A--D (including boundary regularity of $\mathcal{C}$) and the antisymmetry condition~B'. Then there exist constants $\alpha > 0$, $S_{\max} > 0$, and a strictly monotone reparametrization $u = u(g)$ with $u(0) = 0$, such that
\begin{equation}
S_T(g) = S_{\max}\,\tanh\!\left(\alpha\,u(g)\right)\,.
\label{eq:saturation_form}
\end{equation}
The composition law is uniquely determined as the relativistic velocity addition $\mathcal{C}(x,y) = (x+y)/(1+xy)$, up to rescaling by $S_{\max}$. No theory satisfying the axioms can adopt a different composition structure.
\end{theorem}

\subsection{Key Lemmas}

The uniqueness of $\mathrm{arctanh}$ rests on the following chain of results, which we state here and invoke in Step~2 of the proof.

\begin{lemma}[Shift-regularity forces exponential pinning]\label{lem:shift_regular_exponential}
Let $\psi:\mathbb{R}\to(-1,1)$ be a $C^\infty$ strictly increasing odd diffeomorphism and define
\[
C(x,y):=\psi(\psi^{-1}(x)+\psi^{-1}(y))\qquad (x,y\in(-1,1)).
\]
Assume $C\in C^\infty((-1,1)^2)$. Set $\delta(u):=1-\psi(u)$ for $u\to+\infty$.
Then there exists $\kappa>0$ such that for every fixed $v\in\mathbb{R}$,
\[
\lim_{u\to+\infty}\frac{\delta(u+v)}{\delta(u)}=e^{-\kappa v}.
\]
In particular, $\delta(u)=A\,e^{-\kappa u}\,(1+o(1))$ as $u\to+\infty$ for some $A>0$.
\end{lemma}

\begin{proof}[Proof sketch]
Fix $v\in\mathbb{R}$ and consider $x=\psi(u)$ with $u\to+\infty$. Then $C(x,\psi(v))=\psi(u+v)\to 1$. The assumption $C\in C^\infty((-1,1)^2)$ implies that for each fixed $v$, all $x$-derivatives of $x\mapsto C(x,\psi(v))$ remain finite as $x\to 1^-$, which translates into uniform control of the $u$-asymptotics of $\psi(u+v)$ relative to $\psi(u)$. This forces the shift ratio $\delta(u+v)/\delta(u)$ to converge and to satisfy the multiplicative Cauchy equation in $v$, hence it must be $e^{-\kappa v}$ for some $\kappa>0$.
\end{proof}

\begin{proposition}[Quadratic vanishing of the infinitesimal generator]\label{prop:generator_quadratic}
Under the assumptions of Lemma~\ref{lem:shift_regular_exponential}, define the infinitesimal generator $a(x):=\partial_2 C(x,0)$ for $x\in(-1,1)$. Then $a\in C^\infty((-1,1))$, $a(x)>0$ for all $x$, and
\[
a(x)=\lambda(1-x^2)\qquad \text{for some constant }\lambda>0.
\]
Equivalently, writing $\phi:=\psi^{-1}$, one has $\phi'(x)=1/(\lambda(1-x^2))$ for all $x\in(-1,1)$.
\end{proposition}

\begin{proof}[Proof sketch]
The group identity $C(x,0)=x$ gives $a(0)=1$. For each fixed $x$, the map $y\mapsto C(x,y)$ is a $C^\infty$ diffeomorphism (inverse: $y\mapsto C(-x,y)$), hence $a(x)\neq 0$ for all $x$ and by connectedness $a(x)>0$. Writing $x=\psi(u)$, one has $a(\psi(u))=\psi'(u)$. Lemma~\ref{lem:shift_regular_exponential} implies $\psi'(u)\asymp 1-\psi(u)$ as $u\to+\infty$, and by oddness similarly at $u\to-\infty$, yielding $a(x)\asymp 1-x^2$ near $\pm 1$. The $C^\infty$-regularity of $C$ rules out any nonconstant smooth modulation of this quadratic vanishing: if $a(x)=(1-x^2)\,h(x)$ with $h$ nonconstant, then for sufficiently high derivative order $m+n$, the mixed partial $\partial_x^m\partial_y^n C$ diverges as $(x,y)\to(1^-,y_0)$ -- as demonstrated explicitly for the modulated case $h(x)=1+\alpha x^2$ (third derivatives diverge logarithmically). Therefore $h\equiv\lambda$.
\end{proof}

\begin{corollary}[Uniqueness of $\mathrm{arctanh}$ up to scaling]\label{cor:uniqueness_arctanh}
Let $\phi:(-1,1)\to\mathbb{R}$ be odd, $C^\infty$, strictly increasing, $\phi(0)=0$, and $\phi(1^-)=+\infty$. If $C(x,y):=\phi^{-1}(\phi(x)+\phi(y))\in C^\infty((-1,1)^2)$, then there exists $c>0$ such that
\[
\phi(x)=c\,\mathrm{arctanh}(x)\qquad (x\in(-1,1)),
\]
and hence $C(x,y)=(x+y)/(1+xy)$.
\end{corollary}

\begin{proof}
By Proposition~\ref{prop:generator_quadratic}, $\phi'(x)=1/(\lambda(1-x^2))$. Integrating with $\phi(0)=0$ gives $\phi(x)=\lambda^{-1}\,\mathrm{arctanh}(x)$.
\end{proof}

\subsection{Proof of the Saturation Theorem}

The proof proceeds in three steps.

\textit{Step 1: Derivation of the functional equation.}
Define the normalized response $\sigma(g) = S_T(g)/S_{\max} \in [0, 1)$. Axioms~A and~C guarantee that $\sigma$ is bounded and monotonically increasing. Axiom~D [Eq.~\eqref{eq:composition}] requires the existence of a composition:
\begin{equation}
\sigma(g_1 \oplus g_2) = C(\sigma(g_1), \sigma(g_2))\,,
\label{eq:normalized_composition}
\end{equation}
where $C: [0,1) \times [0,1) \to [0,1)$ is commutative, associative, has neutral element $0$, and absorbing boundary $1$.

\textit{Step 2: Reduction to Abel's equation.}
By Axiom~B, $\sigma$ has the local expansion~\eqref{eq:ir_expansion} with $c_1 > 0$. By Axiom~B' [Eq.~\eqref{eq:antisymmetry}], $\sigma(-g) = -\sigma(g)$, so $C$ extends to a group operation on $(-1,1)$ with absorbing boundaries at $\pm 1$.

The associativity condition $C(C(x,y),z) = C(x, C(y,z))$ combined with commutativity and the boundary conditions implies that $C$ defines a \textit{one-parameter continuous group law} on $(-1,1)$ (the Taylor expansion at the origin defines a one-dimensional commutative formal group law in the sense of \cite{Hazewinkel1978}; here the connection is that any smooth, associative, commutative binary operation on a neighborhood of $0 \in \mathbb{R}$ with neutral element $0$ is, by Taylor expansion, a formal group law over $\mathbb{R}$, and conversely, Acz\'el's analytic theorem applies to the convergent realization of this formal structure on the interval $(-1,1)$). We apply Acz\'el's representation theorem \cite{Aczel1966} to the \textit{open interval} $(-1,1)$, where $C$ is smooth, strictly monotone in each argument (since $\sigma$ is strictly monotone by Axiom~C), and the neutral element $0$ lies in the interior. The absorbing boundary conditions $C(x, \pm 1) = \pm 1$ are \textit{not} part of the domain on which Acz\'el's theorem is applied; they are treated as limiting behavior. Specifically, Acz\'el's theorem yields a continuous, strictly monotone function $\phi: (-1,1) \to \mathbb{R}$ with $\phi(0) = 0$. The absorbing boundary conditions then imply $\phi(x) \to +\infty$ as $x \to 1^-$ and $\phi(x) \to -\infty$ as $x \to -1^+$ (since $C(x,y) \to 1$ as $y \to 1$ forces $\phi(x) + \phi(y) \to \infty$, which requires $\phi(y) \to \infty$):
\begin{equation}
\phi(C(x,y)) = \phi(x) + \phi(y)\,.
\label{eq:abel}
\end{equation}
This is Abel's functional equation. Acz\'el's theorem yields $\phi$ continuous; since $C$ is smooth (Axiom~D) and $\phi'(0) = c_1 > 0$ (Axiom~B), differentiating~\eqref{eq:abel} shows that $\phi$ inherits the $C^\infty$ regularity of $C$ on $(-1,1)$. The function $\phi$ is determined up to a positive multiplicative constant. Two conditions jointly select $\phi = \mathrm{arctanh}$:
\begin{enumerate}
\item The antisymmetry (Axiom~B') forces $\phi(-x) = -\phi(x)$.
\item By Corollary~\ref{cor:uniqueness_arctanh}, the $C^\infty$-regularity of $C$ on $(-1,1)^2$ uniquely determines $\phi = c\,\mathrm{arctanh}$ (up to the positive constant $c$). The proof proceeds through the Lemma--Proposition chain of Section~\ref{sec:theorem}: Lemma~\ref{lem:shift_regular_exponential} establishes that $C^\infty$-shift-regularity forces exponential boundary pinning $\delta(u) = A\,e^{-\kappa u}(1+o(1))$; Proposition~\ref{prop:generator_quadratic} shows this implies $\phi'(x) = 1/(\lambda(1-x^2))$, i.e., quadratic vanishing of the infinitesimal generator; integration with $\phi(0)=0$ yields $\phi = c\,\mathrm{arctanh}$. Alternative Abel functions -- algebraic ($x/(1-x^2)$, $\tan(\pi x/2)$: divergent 2nd derivatives at the corner), modulated-logarithmic ($(1+\alpha x^2)\mathrm{arctanh}(x)$: divergent 3rd derivatives) -- are excluded as special cases.
\end{enumerate}
This yields the composition law:
\begin{equation}
C(x, y) = \frac{x + y}{1 + x\,y}\,,
\label{eq:addition_formula}
\end{equation}
which is the relativistic velocity addition formula -- the unique smooth, commutative, associative binary operation on $(-1,1)$ with neutral element $0$, absorbing boundaries $\pm 1$, and $C^\infty$ extension to the closed square.

\textit{Step 3: Inversion to $\tanh$.}
Since $\phi(\sigma(g)) = \mathrm{arctanh}(\sigma(g))$ must be an additive function of the control parameter (i.e., $\phi(\sigma(g)) = \alpha\,u(g)$ for a reparametrization $u$ and constant $\alpha > 0$), we obtain:
\begin{equation}
\sigma(g) = \tanh(\alpha\,u(g))\,,
\end{equation}
and therefore $S_T(g) = S_{\max}\,\tanh(\alpha\,u(g))$, which is~\eqref{eq:saturation_form}. When the axioms hold exactly, the result is exact. Approximate axiom satisfaction (e.g., breakdown of the composition law at very small $g$ where standard GR dominates) would introduce corrections; a formal stability analysis of such perturbations is an open problem. \hfill $\square$

\begin{corollary}[Monotonicity is automatic]
\label{cor:monotonicity}
Axiom~C (monotone cosmic control) follows from Axioms~A, B, B', and~D. That is, any response function satisfying $A + B + B' + D$ is automatically strictly monotonically increasing.
\end{corollary}

\textit{Proof.} Differentiate $\sigma(g + h) = C(\sigma(g), \sigma(h))$ with respect to $h$ at $h = 0$:
\[
\sigma'(g) = \partial_2 C(\sigma(g), 0) \cdot \sigma'(0)\,.
\]
Define $a(x) := \partial_2 C(x, 0)$. By the group structure of Axiom~D, for each fixed $x \in (-1,1)$ the map $y \mapsto C(x,y)$ is a $C^\infty$ diffeomorphism of $(-1,1)$ (its inverse is $y \mapsto C(-x, y)$ from the group inverse). Therefore $\partial_2 C(x, 0) = a(x) \neq 0$ for all $x \in (-1,1)$. Since $(-1,1)$ is connected and $a(0) = \partial_2 C(0,0) = 1 > 0$ (from $C(0,y) = y$), continuity implies $a(x) > 0$ for all $x$. With $\sigma'(0) = c_1 > 0$ (Axiom~B), we obtain $\sigma'(g) = a(\sigma(g)) \cdot c_1 > 0$ for all $g$. \hfill $\square$

\begin{remark}
This means the axiom system can be reduced to three independent axioms (A, B/B', D) plus the derived consequence C. We retain Axiom~C as a separate statement for physical transparency: it makes the monotonicity assumption explicit and connects directly to the thermodynamic arrow of time. But logically, it is a theorem, not an axiom.
\end{remark}

\begin{proposition}[Axiom Independence]
\label{prop:independence}
The axioms are logically independent: for each, there exists a function satisfying the others but violating the chosen one.

\textit{(1) Violates A, satisfies B/B'--D:} $\sigma(g) = g$ (unbounded linear response; $C(x,y) = x+y$, associative and commutative).

\textit{(2) Violates B', satisfies A,C,D:} $\sigma(g) = 1 - e^{-g}$ for $g \geq 0$, with $C(x,y) = x + y - xy$ (probabilistic or); no antisymmetric extension.

\textit{(3) Violates C, satisfies A,B/B' but not D:} $\sigma(g) = \tanh(g) \cdot \cos(\epsilon g)$ for small $\epsilon > 0$; bounded, antisymmetric, smooth with $\sigma'(0) > 0$, but non-monotone for $g \gg 1$. This example shows that C is independent of $\{A, B, B'\}$ alone. However, by Corollary~\ref{cor:monotonicity}, C is \textit{not} independent of $\{A, B, B', D\}$ jointly -- it follows as a derived consequence. Thus C is a dependent axiom, retained for physical clarity.

\textit{(4) Violates D, satisfies A--C,B':} $\sigma(g) = (2/\pi)\arctan(g)$; bounded, monotone, antisymmetric, but its induced composition (via $\phi = \tan(\pi x/2)$) develops divergent second derivatives ($\sim \epsilon^{-1}$) at the boundary of $[-1,1]^2$, failing the $C^\infty$-extension requirement of Axiom~D. Similarly, $\sigma(g) = g/\sqrt{1+g^2}$ (with $\phi(x) = x/\sqrt{1-x^2}$) has divergent second derivatives at the boundary ($\sim \epsilon^{-1/2}$).
\end{proposition}

\subsection{Interpretation}

The theorem states that the $\tanh$ is not a \textit{model choice} but the \textit{canonical solution under boundary regularity} to the combined constraints of boundedness, cooperativity, monotonicity, and scale coherence. The key mathematical insight is that the composition law (Axiom~D) forces the response function to be the inverse of a solution to Abel's equation, and the only such inverse with the correct boundary behavior is $\tanh$.

Physically, the theorem says:
\begin{quote}
\textit{Any quantum gravity theory that is scale-coherent, cooperative, bounded, monotone, and whose composition law extends regularly to the saturation boundary must produce cosmological saturation dynamics of the $\tanh$ type -- regardless of its microscopic mechanism.}
\end{quote}

The reparametrization $u(g)$ absorbs the microscopic details of the theory: different UV completions correspond to different functions $u(g)$, but all produce the same functional form $\tanh(\alpha\,u)$ for the response.


% ===================================================================
% 4. EXCLUSION OF ALTERNATIVES
% ===================================================================
\section{Exclusion of Alternative Profiles}
\label{sec:exclusion}

\begin{corollary}[Exclusion of Alternative Saturation Profiles]
\label{cor:exclusion}
None of the following profiles can simultaneously satisfy Axioms~A--D, except when reducible to the $\tanh$ form via admissible reparametrization $u(g)$:
\begin{enumerate}
\item[(i)] \textbf{Logistic profile:} $S_{\max}/(1 + e^{-\beta u})$.
\item[(ii)] \textbf{Rational profile:} $S_{\max}\,u/(1 + u)$.
\item[(iii)] \textbf{Arctangent profile:} $S_{\max}\,(2/\pi)\arctan(\beta u)$.
\item[(iv)] \textbf{Hard saturation:} $S_{\max}\,\min(u, 1)$.
\end{enumerate}
\end{corollary}

\textit{Proof.} Each alternative fails a specific axiom:

\textit{(i) Logistic.} The logistic function $L(u) = 1/(1 + e^{-\beta u})$ has $L(0) = 1/2 \neq 0$, violating the neutral element condition. The shifted logistic $\tilde{L}(u) = (L(u) - 1/2) \cdot 2$ yields $\tilde{L}(u) = \tanh(\beta u/2)$ -- precisely the $\tanh$ form. Therefore, any ``logistic'' model that satisfies Axiom~D is actually $\tanh$ in disguise.

\textit{(ii) Rational.} The function $R(u) = u/(1+u)$ has Abel transform $\phi(x) = x/(1-x)$, yielding $C(x,y) = \phi^{-1}(\phi(x) + \phi(y)) = (x + y - 2xy)/(1 - xy)$. For small $x, y$, this gives $C(x,y) \approx x + y - 2xy$. The quadratic cross-term $-2xy$ is incompatible with the antisymmetry condition (Axiom~B'): an antisymmetric $\sigma$ implies that the induced composition must satisfy $C(-x,-y) = -C(x,y)$, but $(x+y-2xy)/(1-xy)$ does not. Equivalently, $R(u)$ itself is not antisymmetric, violating Axiom~B'.

\textit{(iii) Arctangent.} The function $A(u) = (2/\pi)\arctan(\beta u)$ has Abel transform $\phi(x) = \tan(\pi x/2)$, yielding $C(x,y) = (2/\pi)\arctan(\tan(\pi x/2) + \tan(\pi y/2))$. While this composition is well-defined and smooth on the open square $(-1,1)^2$, it cannot be extended smoothly to the closed square $[-1,1]^2$: the second partial derivatives $\partial^2 C/\partial x^2$ and $\partial^2 C/\partial x\,\partial y$ both diverge as $(x,y) \to (1,1)$ (numerically $\sim \epsilon^{-1}$ along $x = y = 1 - \epsilon$), violating the $C^\infty$ boundary regularity requirement of Axiom~D. This algebraic ($\epsilon^{-1}$) divergence is characteristic of all Abel functions with power-law boundary behavior; only the logarithmic divergence of $\mathrm{arctanh}$ yields a rational -- hence automatically smooth -- composition.

\textit{(iv) Hard saturation.} $\min(u, 1)$ is not $C^\infty$ at $u = 1$, violating the smoothness requirement of Axiom~B. \hfill $\square$

\begin{remark}
The logistic case (i) is instructive: it appears to be a genuine alternative, but the composition law forces it to reduce to $\tanh$. This illustrates the power of Axiom~D: surface-level freedom in the profile choice is eliminated by the requirement of scale coherence.
\end{remark}


% ===================================================================
% 5. VERIFICATION IN QG PROGRAMS
% ===================================================================
\section{The Axioms in Quantum Gravity Programs}
\label{sec:qg}

We now verify that the four axioms are satisfied -- or at least naturally compatible -- in each major quantum gravity program, when restricted to cosmological observables.

\subsection{Loop Quantum Gravity / Loop Quantum Cosmology}

In Loop Quantum Cosmology (LQC) \cite{Ashtekar2006,Bojowald2008}, the modified Friedmann equation reads:
\begin{equation}
H^2 = \frac{8\pi G}{3}\,\rho\!\left(1 - \frac{\rho}{\rho_c}\right)\,,
\label{eq:lqc_friedmann}
\end{equation}
where $\rho_c \sim \rho_{\mathrm{Pl}}$ is a critical density set by the Barbero-Immirzi parameter $\gamma_{\mathrm{BI}}$. This has exactly the structure of a saturation equation: the expansion rate is bounded ($H^2 \leq H_{\max}^2$), and the bounce at $\rho = \rho_c$ replaces the classical singularity.

\textit{Axiom A:} The critical density $\rho_c$ provides finite geometric capacity. \textit{Axiom B:} For $\rho \ll \rho_c$, the standard Friedmann equation is recovered with $c_1 > 0$. The holonomy corrections are cooperative: the departure from GR grows with curvature. \textit{Axiom C:} In the expanding branch ($\dot{a} > 0$), the solution is monotone. \textit{Axiom D:} LQC does not possess a manifest RG composition law in the sense of Axiom~D; the effective equations inherit diffeomorphism invariance from full LQG, which provides structural coherence across scales, but a rigorous derivation of the associative composition property~\eqref{eq:composition} from the LQG Hilbert space remains an open problem (hence the ``?'' in Table~\ref{tab:axiom_check}).

Defining $X = 1 - 2\rho/\rho_c$, the LQC Friedmann equation~\eqref{eq:lqc_friedmann} can be rewritten as $H^2 \propto (1 - X^2)$. The factor $(1-X^2)$ is the same nonlinearity that appears on the right-hand side of the saturation ODE~\eqref{eq:saturation_ode}; note, however, that the LQC equation governs $H^2$ while the CRM ODE governs $dX/da$ -- the structural analogy is in the saturation mechanism, not a direct identification of the two equations. The CRM parameters $k$ and $\Phi_0$ would map to the LQG minimum area gap $\Delta$ (the smallest nonzero eigenvalue of the area operator) and the Barbero-Immirzi parameter $\gamma_{\mathrm{BI}}$, though this mapping has not been rigorously derived.

\subsection{Asymptotic Safety}

In the asymptotic safety program \cite{Weinberg1979,Reuter1998,Percacci2017}, the gravitational effective action has a UV fixed point $g^* > 0$ where the dimensionless Newton's constant and cosmological constant become scale-independent. The RG flow from UV to IR defines the effective gravitational coupling $G_k$ and cosmological ``constant'' $\Lambda_k$ as functions of the RG scale $k$.

\textit{Axiom A:} The UV fixed point provides a finite upper bound on the effective coupling. \textit{Axiom B:} The flow away from the fixed point is cooperative: the relevant direction has positive critical exponent \cite{Percacci2017}. \textit{Axiom C:} The flow is monotone along the relevant direction (by definition of relevance). \textit{Axiom D:} A clarification is needed. The Wilson RG semigroup $R_{k_1} \circ R_{k_2} = R_{k_1 k_2}$ operates on the space of couplings (the ``theory space''), not directly on the response function $S_T$. To connect the two, one must project the RG flow onto the single relevant direction parametrized by $g$. In the asymptotic safety scenario with a UV fixed point, the flow near the fixed point is controlled by a finite number of relevant directions; in the simplest truncation (Einstein-Hilbert with $G_k, \Lambda_k$), there are typically two relevant directions. The single-field assumption of Axiom~D corresponds to restricting attention to the dominant relevant direction, under which the RG semigroup on theory space \textit{induces} a composition law on the one-dimensional response $S_T(g)$. This projection is exact when a single relevant direction dominates (as in the cosmological sector), and approximate otherwise -- a limitation shared with the single-field restriction of the theorem (see Limitation~4).

Indeed, asymptotic safety is the QG program where Axiom~D is most naturally satisfied, because it \textit{is} a renormalization group statement. The asymptotic safety scenario for cosmology \cite{Bonanno2002} produces effective Friedmann equations with saturation behavior at high curvature, consistent with the $\tanh$ form.

Recent computations within increasingly sophisticated truncations provide growing numerical evidence for the Reuter fixed point. Eichhorn and Held~\cite{EichhornHeld2018} demonstrate that asymptotic safety constrains the top-quark Yukawa coupling, yielding a prediction of approximately 171\,GeV for the top-quark pole mass, consistent with the observed value of $172.6 \pm 0.3$\,GeV \cite{PDG2024} -- a concrete particle-physics consequence of the UV fixed point. More broadly, Eichhorn~\cite{Eichhorn2020} reviews how the fixed-point structure naturally accommodates the Standard Model matter content without additional fine-tuning. These results strengthen the case that the asymptotic safety program naturally realizes the axioms: the fixed-point existence (Axiom~A), cooperative flow (Axiom~B), and RG composition (Axiom~D) are properties verified in state-of-the-art functional RG computations.

\subsection{String Theory and Holography}

In string theory, the UV finiteness of the string amplitude provides Axiom~A. The string landscape, while introducing a selection problem, does not violate the axioms for any \textit{specific} vacuum: within a given compactification, the low-energy EFT has a unique RG flow satisfying Axioms~B--D.

The holographic approach (AdS/CFT \cite{Maldacena1999}) provides particularly strong support for Axiom~D: the bulk/boundary correspondence \textit{is} a composition law relating geometric descriptions at different scales. The Ryu-Takayanagi formula \cite{Ryu2006} and its quantum corrections establish a monotone relationship between entanglement entropy and geometric area, satisfying Axiom~C.

For cosmological applications, the de~Sitter entropy bound $S_{\mathrm{dS}} = \pi/(\Lambda G) < \infty$ provides Axiom~A, while the holographic RG flow \cite{deBoer2000} provides the composition structure.

\subsection{Causal Set Theory}

In causal set theory \cite{Sorkin2003,Surya2019}, the cosmological constant arises as a dynamical quantity $\Lambda \sim 1/\sqrt{N}$, where $N$ is the number of causal set elements. This produces a cosmological ``constant'' that decreases as the universe grows, approaching a finite asymptotic value -- precisely the saturation behavior of Axiom~C.

\textit{Axiom A:} Finite density of causal set elements provides geometric capacity. \textit{Axiom B:} While the Poisson sprinkling itself is uncorrelated, the \textit{dynamics} of causal set evolution (sequential growth models \cite{Surya2019}) produce cooperative ordering: newly added elements preferentially attach to the existing causal structure, enhancing geometric order. \textit{Axiom C:} The $1/\sqrt{N}$ scaling is monotonically decreasing. \textit{Axiom D:} The sequential growth dynamics provide a natural composition law, though its precise relationship to the RG composition of Axiom~D requires further analysis.

\subsection{Causal Dynamical Triangulations}

Causal Dynamical Triangulations (CDT) \cite{Ambjorn2005,Loll2019} construct quantum gravity via a regularized path integral over causal triangulations. Monte Carlo simulations reveal a phase diagram with a physically relevant ``de~Sitter phase'' (phase~$C_{dS}$) in which the emergent geometry has Hausdorff dimension $d_H \approx 4$ and a volume profile well-described by Euclidean de~Sitter space.

\textit{Axiom A:} The triangulation provides a finite number of building blocks per spacetime volume, enforcing finite geometric capacity. \textit{Axiom B:} In the de~Sitter phase, the volume profile grows cooperatively before saturating -- the geometric ordering is self-reinforcing. The odd-power structure (Axiom~B') has not been explicitly verified. \textit{Axiom C:} The de~Sitter volume profile is monotone (in the expanding branch). \textit{Axiom D:} CDT possesses natural blocking transformations (coarse-graining of triangulations), providing a candidate for the composition law, but the precise associative structure has not been rigorously established.

\subsection{Noncommutative Geometry}

Connes' program \cite{Connes1994} reconstructs geometry from spectral data (the Dirac operator on a noncommutative algebra). The spectral action principle yields the Standard Model Lagrangian coupled to gravity from the spectrum of a single operator.

\textit{Axiom A:} The spectral action has a natural UV cutoff $\Lambda$ (the energy scale at which the spectral sum is evaluated), providing finite geometric capacity. \textit{Axiom B:} The Seeley-DeWitt coefficients $a_0, a_2, a_4, \ldots$ of the heat kernel expansion have even powers of $\Lambda$; the odd-power structure required by Axiom~B emerges after reparametrization to a suitable coupling variable, but this identification requires further analysis. \textit{Axiom C:} The spectral flow is monotone in the appropriate direction. \textit{Axiom D:} The composition of spectral triples provides algebraic composition coherence.

\subsection{Summary: Why Universality?}

Table~\ref{tab:axiom_check} summarizes the axiom verification. No single QG program has been rigorously proven to satisfy all four axioms simultaneously; rather, each program naturally satisfies a substantial subset, with the remaining axioms being plausible but not yet formally established (indicated by ($\checkmark$) and ? in the table). The theorem's value is therefore twofold: (i)~as a \textit{conditional} result, it identifies the precise conditions under which $\tanh$ saturation is inevitable, and (ii)~as a \textit{diagnostic}, it pinpoints which properties of each QG program require further investigation. The broad compatibility across programs explains the ``unreasonable universality'' of saturation dynamics: it is not that the programs agree on microscopic details, but that they all satisfy -- or plausibly satisfy -- the same macroscopic constraints.

\begin{table}[h]
\caption{Axiom verification across quantum gravity programs. $\checkmark$ = naturally satisfied with explicit construction, ($\checkmark$) = plausible with mild assumptions, ? = open question requiring further analysis.}
\label{tab:axiom_check}
\begin{ruledtabular}
\begin{tabular}{lcccc}
Program & A & B/B' & C & D \\
\hline
LQG/LQC & $\checkmark$ & $\checkmark$ & $\checkmark$ & ? \\
Asymptotic safety & $\checkmark$ & ($\checkmark$) & $\checkmark$ & $\checkmark$ \\
Strings/holography & $\checkmark$ & ($\checkmark$) & ($\checkmark$) & ($\checkmark$) \\
Causal sets & $\checkmark$ & ? & $\checkmark$ & ? \\
CDT & $\checkmark$ & ($\checkmark$) & ($\checkmark$) & ? \\
NCG & $\checkmark$ & ? & ? & ? \\
CRM (effective) & $\checkmark$ & $\checkmark$ & $\checkmark$ & $\checkmark$ \\
\end{tabular}
\end{ruledtabular}
\end{table}


% ===================================================================
% 6. ONTOLOGY: ORDER, NOT METRIC
% ===================================================================
\section{Spacetime as an Ordered Phase}
\label{sec:ontology}

The broad axiom satisfaction documented in Table~\ref{tab:axiom_check} raises a conceptual question: if the $\tanh$ form is universal, what does this universality \textit{mean} for the nature of spacetime?

\subsection{The conceptual shift}

The Saturation Theorem, combined with the CRM's empirical success, suggests a reconceptualization of quantum gravity:

\begin{quote}
\textit{Quantum gravity is not the quantization of the metric, but the theory of how geometric order emerges from a pre-geometric quantum system.}
\end{quote}

The $\tanh$ saturation is not a property of a specific field or coupling, but the universal signature of a capacity-limited cooperative ordering process -- independent of whether the microscopic degrees of freedom are spin networks (LQG), strings, causal set elements, or code qubits (QEC).

\subsection{The ferromagnetic analogy}

The analogy with spontaneous magnetization is precise at the level of the Ginzburg-Landau free energy. The CRM saturation ODE $dX/da = k(1 - X^2)$ is the equation of motion for the order parameter of a second-order phase transition with double-well free energy $F(X) = -k(X - X^3/3)$. Paper~III \cite{Geiger2026c} develops this analogy in detail, identifying:

\begin{center}
\begin{tabular}{ll}
\hline
Ferromagnetism & CRM Cosmology \\
\hline
Magnetization $M$ & Curvature return $\Omega_\Phi$ \\
Inverse temperature $J/k_BT$ & Scale factor $a$ \\
Curie point ($J/k_BT_c$) & Transition $a_{\mathrm{trans}}$ \\
Spin interaction $J$ & Curvature coupling $k$ \\
Saturation $M_s$ & Saturation $\Phi_0$ \\
$\tanh(J/k_BT)$ & $\tanh(k(a - a_{\mathrm{trans}}))$ \\
\hline
\end{tabular}
\end{center}

The Saturation Theorem provides the mathematical foundation for this analogy at the \textit{mean-field level}: both systems satisfy Axioms~A--D in the mean-field approximation, and therefore both produce $\tanh$ response functions. We note that the exact magnetization of a $d$-dimensional Ising model deviates from the mean-field $\tanh$ near the critical temperature ($M \sim (T_c - T)^{\beta}$ with non-trivial critical exponent $\beta \neq 1/2$ for $d < 4$). This deviation corresponds to the breakdown of Axiom~D's $C^\infty$ boundary regularity near the \textit{critical point} (the onset of ordering), not near \textit{saturation} ($T \to 0$). The analogy is therefore precise for the saturation regime and mean-field theory, while fluctuation corrections near criticality require a multi-field extension (see Outlook). The universality is structural, not metaphorical, within the domain of validity of the axioms.

Having established the ontological framework, we now turn to its empirical consequences.

% ===================================================================
% 7. OBSERVATIONAL CONSEQUENCES
% ===================================================================
\section{Observational Consequences and Open Questions}
\label{sec:observational}

\subsection{What the theorem predicts}

The Saturation Theorem has two classes of predictions:

\textbf{Universal predictions} (independent of UV completion):
\begin{enumerate}
\item The cosmological expansion history must follow $\tanh$-type saturation. We emphasize that the CRM's $\Delta\chi^2 = -5.5$ result \cite{Geiger2026b} does \textit{not} constitute an independent test of the theorem, since the CRM was constructed with a $\tanh$ profile before the axioms were formulated. The theorem's empirical content is the \textit{composition law} $\mathcal{C}(x,y) = (x+y)/(1+xy)$, which is a prediction beyond any individual profile fit. An independent test would require either (a)~fitting a model-agnostic $\tanh$-template to SN+CMB+BAO data without CRM-specific priors, or (b)~verifying the composition law directly by comparing saturation responses at different redshift bins and checking consistency with the velocity-addition rule.
\item The running curvature coupling $\beta_{\mathrm{eff}}(a)$ must be a monotone transition function (already confirmed: $\beta_{\mathrm{early}} \approx 2.8 \to \beta_{\mathrm{late}} \approx 2.0$ \cite{Geiger2026b}).
\item No alternative relaxation profile (logistic, power-law, etc.) can fit the data \textit{and} be scale-coherent simultaneously.
\end{enumerate}

\textbf{UV-dependent predictions} (discriminating between completions):
\begin{enumerate}
\item The specific values of $k$ (saturation rate), $\Phi_0$ (saturation amplitude), and $\gamma$ (coupling constant) depend on the UV theory.
\item The transition sharpness ($n$ in the sigmoidal transition) encodes the universality class.
\item Sub-Planckian corrections to the $\tanh$ profile are UV-theory-specific.
\end{enumerate}

\subsection{The missing piece: why \textit{these} parameters?}

The Saturation Theorem determines the \textit{form} but not the \textit{coefficients}. The remaining open question is:

\begin{quote}
\textit{Why does the universe have the specific values $k \approx 9.8$, $\Phi_0 \approx 0.43$, $\gamma \sim \mathcal{O}(1)\,H_0^{-2}$?}
\end{quote}

This question is analogous to the fine-structure constant problem: the \textit{existence} of electromagnetic coupling is forced by gauge symmetry, but the \textit{value} $\alpha \approx 1/137$ is not determined by the symmetry alone.

Several UV completions make specific predictions:
\begin{itemize}
\item \textbf{LQC:} $\Phi_0$ maps to the Barbero-Immirzi parameter $\gamma_{\mathrm{BI}}$, which is constrained by black hole entropy \cite{Meissner2004}. This would predict a specific value of $\Phi_0$.
\item \textbf{Asymptotic safety:} The critical exponents at the UV fixed point determine the transition sharpness $n$, providing a specific prediction for the shape of $\beta_{\mathrm{eff}}(a)$ \cite{Reuter1998}.
\item \textbf{Holography:} The de~Sitter entropy $S_{\mathrm{dS}} = \pi/(\Lambda G)$ relates $\Phi_0$ to the ratio of UV and IR scales.
\end{itemize}

A concrete test case is the $\sqrt{\pi}$ conjecture of Paper~II \cite{Geiger2026b}: the CRM predicts a MOND-like enhancement factor $\mu_{\mathrm{eff}} = \sqrt{\pi}$ arising from Gaussian normalization of the saturation partition function. Deriving this value from a specific UV completion would directly connect quantum gravity parameters to observed galactic dynamics.

\subsection{Scalar field oscillations as a QG signature}

Paper~III \cite{Geiger2026c} shows that the P\"oschl-Teller potential supports a discrete spectrum of bound states. In the cosmological context, these correspond to oscillatory corrections to the expansion rate at late times:
\begin{equation}
H^2(a) = H_{\mathrm{smooth}}^2(a)\left[1 + \epsilon \cdot e^{-\Gamma a} \cos(\omega a + \delta)\right]\,,
\end{equation}
with $\epsilon \ll 1$. These oscillations would be a direct signature of the \textit{quantum nature} of the saturation mechanism -- analogous to the discrete energy levels of a quantum system. Detection in high-precision BAO or SN data would constitute evidence for the quantum origin of the CRM saturation.


% ===================================================================
% 8. DISCUSSION
% ===================================================================
\section{Discussion and Outlook}
\label{sec:discussion}

\subsection{What has been achieved}

\begin{enumerate}
\item \textbf{A No-Go theorem:} Any cosmological theory satisfying Axioms~A--D with boundary regularity must produce $\tanh$-type saturation. This is a structural result, independent of specific QG programs.

\item \textbf{Universality explained:} The observation that LQG, asymptotic safety, strings, causal sets, and NCG all produce structurally similar saturation dynamics is no longer mysterious -- it follows from the Saturation Theorem.

\item \textbf{CRM elevated:} The CRM's saturation ODE $dX/da = k(1-X^2)$ is not an ad hoc phenomenological fit, but the universal normal form of capacity-limited cooperative relaxation -- the canonical smooth, scale-coherent saturation profile under boundary regularity.

\item \textbf{Ontological clarity:} Quantum gravity is not the quantization of the metric, but the theory of geometric order -- a phase transition from a pre-geometric quantum system to classical spacetime.
\end{enumerate}

\subsection{Relation to the CRM program}

This paper completes the conceptual arc of the CRM series: game-theoretic foundation (Paper~I \cite{Geiger2026}), phenomenological fit with $\Delta\chi^2 = -5.5$ (Paper~II \cite{Geiger2026b}), Lagrangian derivation (Paper~III \cite{Geiger2026c}), galactic dynamics (Paper~IV \cite{Geiger2026d}), and now axiomatic foundation (this work). The logical status of the CRM is:
\begin{enumerate}
\item The \textit{form} of the dynamics ($\tanh$) is a mathematical theorem (this paper).
\item The \textit{Lagrangian} is derived (Paper~III).
\item The \textit{parameters} are constrained by data (Papers~II, III, IV).
\item The \textit{UV completion} remains open -- but the theorem shows that \textit{any} viable completion must reproduce the CRM effective dynamics.
\end{enumerate}

\subsection{Limitations and caveats}

\begin{enumerate}
\item \textbf{Axiom D is strong.} The renormalization-group composition law is the most restrictive axiom. Theories without a well-defined RG flow (some discrete approaches, multiverse scenarios) may not satisfy it. We regard this as a feature, not a bug: Axiom~D is precisely the condition that distinguishes predictive theories from frameworks without observable consequences.

\item \textbf{The proof is algebraic, not constructive.} The theorem establishes the form of the response function but does not construct the UV theory. The ``reparametrization'' $u(g)$ absorbs all microscopic information; determining $u(g)$ from first principles requires solving the specific QG theory. A potential objection is that any bounded monotone function can trivially be written as $\tanh(\alpha\,u(g))$ for a suitable $u(g)$, making the $\tanh$ form contentless. We address this directly: the theorem's non-trivial, falsifiable content resides in the \textit{composition law} $\mathcal{C}(x,y) = (x+y)/(1+xy)$, not in the profile shape. The composition law is a structural constraint on how responses at different scales combine -- it is the gravitational analog of the relativistic velocity addition rule. This constraint (i)~is not satisfied by generic bounded monotone functions (e.g., $\text{erf}(g)$ does not admit an associative smooth composition with absorbing boundary), (ii)~is testable in principle by comparing saturation responses at different scales, and (iii)~restricts the space of allowed effective theories to those whose RG flow has the one-dimensional group structure identified by the theorem.

\item \textbf{Cosmological restriction.} The axioms are formulated for cosmological observables. Whether they extend to local physics (black holes, gravitational waves) requires further analysis. The CRM's existing predictions for local physics (chameleon screening, $\alpha_T = 0$, Bullet Cluster resolution \cite{Geiger2026b,Geiger2026c,Geiger2026d}) are derived from the Lagrangian, not from the axioms.

\item \textbf{The continuous group classification and single-field restriction.} The proof relies on Acz\'el's representation theorem for continuous associative operations \cite{Aczel1966} and the classification of one-dimensional formal group laws \cite{Hazewinkel1978}. While these results are well-established, the physical applicability requires the assumption that the composition is \textit{one-dimensional} -- i.e., that there is a single relevant direction. This is a genuine limitation: the CRM itself (Paper~IV \cite{Geiger2026d}) introduces a vector sector alongside the scalar, and realistic modified gravity theories (e.g., Horndeski, beyond-Horndeski) generically have multiple coupled degrees of freedom. However, we note that the theorem applies to the \textit{scalar sector in isolation}: the CRM's scalar saturation channel $\Omega_\Phi(a) = \Phi_0 \tanh(k(a - a_{\mathrm{trans}}))$ is a one-dimensional response that satisfies the axioms independently of the vector sector. The multi-field extension (classifying multi-dimensional commutative group laws, including elliptic and Lubin-Tate formal groups) is an open mathematical problem that would strengthen the theorem but is not required for the scalar-sector result.

\item \textbf{The uniqueness selection.} The selection of $\mathrm{arctanh}$ as the unique odd Abel function whose induced composition extends to $C^\infty([-1,1]^2)$ rests on four pillars: (i)~explicit verification against all standard alternatives (algebraic: $x/(1-x^2)$, $\tan(\pi x/2)$, $x/\sqrt{1-x^2}$; modulated logarithmic: $(1+\alpha x^2)\,\mathrm{arctanh}(x)$); (ii)~the structural argument that only the logarithmic divergence rate of $\mathrm{arctanh}$ produces a \textit{rational} composition $(x+y)/(1+xy)$, hence automatic $C^\infty$ extension; (iii)~the perturbative argument (Step~2) that any deviation from $\mathrm{arctanh}$ flat at the boundary introduces logarithmic divergences in higher derivatives; and (iv)~the \textit{shift-regularity argument} (Step~2): $C^\infty$-regularity of the translation family $T_v(u) = \psi(u+v)$ at the boundary forces the defect $\delta(u) = 1 - \psi(u)$ to decay exponentially, which integrates to $\phi = c\,\mathrm{arctanh}$ globally. Pillar~(iv) provides the strongest argument, reducing the uniqueness to a single analytic principle (exponential boundary attachment from shift regularity).
\end{enumerate}

\subsection{The No-Go theorem as a filter}

The Saturation Theorem serves as a \textit{filter} for quantum gravity proposals:

\begin{itemize}
\item Any proposed QG theory that violates Axiom~A (allows infinite geometric order) must address the singularity problem separately.
\item Any theory that violates Axiom~C (has oscillating or multi-fixed-point dynamics) must explain why the observed universe has a unique thermodynamic arrow.
\item Any theory that violates Axiom~D (lacks scale coherence) must explain why its predictions are insensitive to the observation scale.
\end{itemize}

Theories satisfying all four axioms are \textit{constrained} to produce $\tanh$ saturation, providing a concrete, falsifiable prediction for their cosmological sector.

\subsection{Relation to other frameworks}

\textit{Swampland conjectures.} The de~Sitter swampland conjecture \cite{Obied2018} posits $|\nabla V|/V \geq c \sim \mathcal{O}(1)$ in any quantum gravity landscape, disfavoring a strict cosmological constant. The Saturation Theorem is complementary: it does not address the existence of de~Sitter vacua, but constrains the \textit{dynamical approach} to the late-time attractor. A CRM-type saturation with $\Phi_0 < \infty$ is compatible with the swampland bound, since the effective dark energy density $\rho_\Phi(a)$ evolves dynamically rather than being a fixed $\Lambda$.

\textit{EFT of Dark Energy.} The effective field theory of dark energy \cite{Gubitosi2013} parametrizes deviations from $\Lambda$CDM at the level of the perturbed action. The Saturation Theorem operates at a different level: it constrains the \textit{background} response function, not the perturbative operator basis. A future task is to express the $\tanh$ saturation in the EFT-of-DE language, mapping the CRM parameters $(k, \Phi_0, \gamma)$ to the EFT operator coefficients $\{\alpha_i(t)\}$.

\subsection{Outlook}

Three directions are immediate:

\begin{enumerate}
\item \textbf{Parameter derivation.} The next step is to derive $k$, $\Phi_0$, and $\gamma$ from a specific UV completion. The most promising candidate is the LQG/holographic intersection, where the Barbero-Immirzi parameter and the holographic entropy bound may jointly fix the saturation parameters.

\item \textbf{Multi-field extension.} Extending the theorem to theories with multiple saturation channels (e.g., scalar + vector in the CRM \cite{Geiger2026d}) requires the classification of multi-dimensional commutative group laws (the higher-dimensional analog of the Abel-equation reduction). This is mathematically richer and may yield constraints on the vector sector parameters $K_B$ and $\mathcal{F}$.

\item \textbf{Black hole interior.} If the Saturation Theorem applies to black hole interiors (with the radial coordinate replacing the scale factor), it predicts $\tanh$-type singularity resolution -- consistent with LQG bounce models \cite{Ashtekar2006} and the Planck star scenario \cite{Rovelli2014}.
\end{enumerate}


% ===================================================================
% 9. CONCLUSION
% ===================================================================
\section{Conclusion}
\label{sec:conclusion}

We have proven the Saturation Theorem: any quantum gravity theory satisfying four minimal axioms plus boundary regularity of the composition -- finite geometric capacity, cooperative reinforcement, monotone cosmic control, and renormalization-group composition -- must produce cosmological dynamics with a $\tanh$-type saturation profile. This is not a model choice but a mathematical necessity, analogous to how gauge symmetry forces the existence of gauge bosons without specifying the coupling constant.

The theorem has three consequences:

\begin{enumerate}
\item The CRM's saturation ODE $dX/da = k(1-X^2)$ is the \textit{universal normal form} of capacity-limited cooperative relaxation -- the unique response function compatible with boundedness, cooperativity, monotonicity, and scale coherence.

\item The observation that LQG, asymptotic safety, string theory, causal sets, and noncommutative geometry all produce $\tanh$-type saturation in their cosmological sectors is \textit{explained}: it follows from the axioms that each program satisfies.

\item Quantum gravity is reconceptualized as the theory of geometric order -- a phase transition from a pre-geometric quantum system to classical spacetime -- rather than the ``quantization of the metric.''
\end{enumerate}

The remaining open problem -- deriving the specific parameters $(k, \Phi_0, \gamma)$ from a UV completion -- is the analog of deriving the fine-structure constant from a theory of quantum electrodynamics. The theorem constrains the \textit{form} of any viable theory; the community's task is now to determine which UV completion selects the \textit{observed} values.

\begin{quote}
\textit{The saturation profile is constrained by the axioms to be of the $\tanh$ type -- not as a model choice, but as a consequence of scale coherence and boundary regularity.}
\end{quote}


% ===================================================================
% ACKNOWLEDGMENTS
% ===================================================================
\begin{acknowledgments}
This work was conducted as independent research without institutional funding. The author thanks the open-source communities behind the scientific computing infrastructure used throughout the CRM program.

\textit{AI tools.} AI-based tools (Claude by Anthropic; Copilot by Microsoft/GitHub; Gemini by Google DeepMind) were used for mathematical formalization, literature survey, code development, and text generation. All physical hypotheses, axiomatic choices, scientific interpretation, and responsibility for the content lie solely with the author.

\textit{Funding information.} This research received no external funding.

\textit{Data availability.} This paper is purely theoretical. The CRM analysis code and data are available at \url{https://github.com/lukisch/crm-cosmology} (DOI: \href{https://doi.org/10.5281/zenodo.18728935}{10.5281/zenodo.18728935}).
\end{acknowledgments}


% ===================================================================
% BIBLIOGRAPHY
% ===================================================================
\begin{thebibliography}{50}

\bibitem{Geiger2026}
L.~Geiger (2026).
Game-Theoretic Cosmology and the Curvature Relaxation Model: Nash Equilibria Between Null Space and Spacetime Bubble.
Companion paper (Paper~I).
\url{https://github.com/lukisch/crm-cosmology}.
\href{https://doi.org/10.5281/zenodo.18728935}{doi:10.5281/zenodo.18728935}.

\bibitem{Geiger2026b}
L.~Geiger (2026).
Eliminating the Dark Sector: Unifying the Curvature Relaxation Model with MOND.
Companion paper (Paper~II).
\href{https://doi.org/10.5281/zenodo.18728935}{doi:10.5281/zenodo.18728935}.

\bibitem{Geiger2026c}
L.~Geiger (2026).
Microscopic Foundations of the Curvature Relaxation Model: From Quantum Geometry to Macroscopic Saturation.
Companion paper (Paper~III).
\href{https://doi.org/10.5281/zenodo.18728935}{doi:10.5281/zenodo.18728935}.

\bibitem{Geiger2026d}
L.~Geiger (2026).
The Galactic--Cosmological Nexus: Connecting CRM to MOND Dynamics via Curvature Saturation.
Companion paper (Paper~IV).
\href{https://doi.org/10.5281/zenodo.18728935}{doi:10.5281/zenodo.18728935}.

\bibitem{Donoghue1994}
J.~F.~Donoghue (1994).
General relativity as an effective field theory: The leading quantum corrections.
\textit{Physical Review D}, 50(6), 3874.
DOI: 10.1103/PhysRevD.50.3874.
arXiv:gr-qc/9405057.

\bibitem{Eichhorn2020}
A.~Eichhorn,
\textit{Asymptotically safe gravity},
arXiv:2003.00044 (2020).

\bibitem{EichhornHeld2018}
A.~Eichhorn and A.~Held,
\textit{Top mass from asymptotic safety},
Phys.\ Lett.\ B \textbf{777}, 217--221 (2018).
arXiv:1707.01107.

\bibitem{Weinberg1979}
S.~Weinberg (1979).
Ultraviolet divergences in quantum theories of gravitation.
In S.~W.~Hawking \& W.~Israel (Eds.), \textit{General Relativity: An Einstein Centenary Survey}, pp.\ 790--831.
Cambridge University Press.

\bibitem{Reuter1998}
M.~Reuter (1998).
Nonperturbative evolution equation for quantum gravity.
\textit{Physical Review D}, 57(2), 971.
DOI: 10.1103/PhysRevD.57.971.
arXiv:hep-th/9605030.

\bibitem{Percacci2017}
R.~Percacci (2017).
\textit{An Introduction to Covariant Quantum Gravity and Asymptotic Safety}.
World Scientific.

\bibitem{Polchinski1998}
J.~Polchinski (1998).
\textit{String Theory}, Vols.~1--2.
Cambridge University Press.

\bibitem{Rovelli2004}
C.~Rovelli (2004).
\textit{Quantum Gravity}.
Cambridge University Press.

\bibitem{Thiemann2007}
T.~Thiemann (2007).
\textit{Modern Canonical Quantum General Relativity}.
Cambridge University Press.

\bibitem{Maldacena1999}
J.~Maldacena (1998).
The large-$N$ limit of superconformal field theories and supergravity.
\textit{Adv.\ Theor.\ Math.\ Phys.}\ \textbf{2}, 231--252.
[Reprinted in \textit{Int.\ J.\ Theor.\ Phys.}\ \textbf{38}, 1113--1133 (1999).]
arXiv:hep-th/9711200.

\bibitem{Sorkin2003}
R.~D.~Sorkin (2003).
Causal sets: Discrete gravity.
In \textit{Lectures on Quantum Gravity}, pp.\ 305--327.
Springer.
arXiv:gr-qc/0309009.

\bibitem{Surya2019}
S.~Surya (2019).
The causal set approach to quantum gravity.
\textit{Living Reviews in Relativity}, 22, 5.
DOI: 10.1007/s41114-019-0023-1.

\bibitem{Connes1994}
A.~Connes (1994).
\textit{Noncommutative Geometry}.
Academic Press.

\bibitem{Bekenstein1981}
J.~D.~Bekenstein (1981).
A universal upper bound on the entropy-to-energy ratio for bounded systems.
\textit{Physical Review D}, 23(2), 287.
DOI: 10.1103/PhysRevD.23.287.

\bibitem{Almheiri2015}
A.~Almheiri, X.~Dong, \& D.~Harlow (2015).
Bulk locality and quantum error correction in AdS/CFT.
\textit{Journal of High Energy Physics}, 2015, 163.
DOI: 10.1007/JHEP04(2015)163.

\bibitem{Aczel1966}
J.~Acz\'el (1966).
\textit{Lectures on Functional Equations and Their Applications}.
Academic Press, New York.

\bibitem{Hazewinkel1978}
M.~Hazewinkel (1978).
\textit{Formal Groups and Applications}.
Academic Press.

\bibitem{Ashtekar2006}
A.~Ashtekar, T.~Pawlowski, \& P.~Singh (2006).
Quantum nature of the Big Bang.
\textit{Physical Review Letters}, 96(14), 141301.
DOI: 10.1103/PhysRevLett.96.141301.
arXiv:gr-qc/0602086.

\bibitem{Bojowald2008}
M.~Bojowald (2008).
Loop quantum cosmology.
\textit{Living Reviews in Relativity}, 11, 4.
DOI: 10.12942/lrr-2008-4.

\bibitem{Bonanno2002}
A.~Bonanno \& M.~Reuter (2002).
Cosmology of the Planck era from a renormalization group for quantum gravity.
\textit{Physical Review D}, 65(4), 043508.
DOI: 10.1103/PhysRevD.65.043508.
arXiv:hep-th/0106133.

\bibitem{Ryu2006}
S.~Ryu \& T.~Takayanagi (2006).
Holographic derivation of entanglement entropy from the anti-de~Sitter space/conformal field theory correspondence.
\textit{Physical Review Letters}, 96(18), 181602.
DOI: 10.1103/PhysRevLett.96.181602.
arXiv:hep-th/0603001.

\bibitem{deBoer2000}
J.~de Boer, E.~P.~Verlinde, \& H.~L.~Verlinde (2000).
On the holographic renormalization group.
\textit{Journal of High Energy Physics}, 2000(08), 003.
DOI: 10.1088/1126-6708/2000/08/003.
arXiv:hep-th/9912012.


\bibitem{Meissner2004}
K.~A.~Meissner (2004).
Black-hole entropy in Loop Quantum Gravity.
\textit{Classical and Quantum Gravity}, 21(22), 5245.
DOI: 10.1088/0264-9381/21/22/015.
arXiv:gr-qc/0407052.

\bibitem{Rovelli2014}
C.~Rovelli \& F.~Vidotto (2014).
Planck stars.
\textit{International Journal of Modern Physics D}, 23(12), 1442026.
DOI: 10.1142/S0218271814420267.
arXiv:1401.6562.

\bibitem{PDG2024}
S.~Navas et~al. (Particle Data Group) (2024).
Review of Particle Physics.
\textit{Physical Review D}, 110, 030001.
DOI: 10.1103/PhysRevD.110.030001.

\bibitem{Obied2018}
G.~Obied, H.~Ooguri, L.~Spodyneiko, \& C.~Vafa (2018).
De Sitter Space and the Swampland.
arXiv:1806.08362.

\bibitem{Gubitosi2013}
G.~Gubitosi, F.~Piazza, \& F.~Vernizzi (2013).
The Effective Field Theory of Dark Energy.
\textit{Journal of Cosmology and Astroparticle Physics}, 2013(02), 032.
DOI: 10.1088/1475-7516/2013/02/032.
arXiv:1210.0201.

\bibitem{Ambjorn2005}
J.~Ambj\o{}rn, J.~Jurkiewicz, \& R.~Loll (2005).
The spectral dimension of the universe is scale dependent.
\textit{Physical Review Letters}, 95(17), 171301.
DOI: 10.1103/PhysRevLett.95.171301.
arXiv:hep-th/0505113.

\bibitem{Loll2019}
R.~Loll (2020).
Quantum gravity from causal dynamical triangulations: a review.
\textit{Classical and Quantum Gravity}, 37(1), 013002.
DOI: 10.1088/1361-6382/ab57c7.
arXiv:1905.08669.

\end{thebibliography}

\end{document}
